% reference_exam.tex
% Structural reference for the lecture-materials-assistant exam artifact.
% All course-specific content replaced with bracketed placeholders.
%
% ANSWER KEY: change \answersfalse to \answerstrue and recompile.
%
% Compile with: pdflatex reference_exam.tex

\documentclass{article}
\usepackage[letterpaper, margin=1in]{geometry}
\usepackage{enumitem}
\usepackage{listings}
\usepackage{fontenc}

% --- Answer key toggle ---
\newif\ifanswers
\answersfalse   % <-- change to \answerstrue for answer key

% --- Code listing style ---
\lstset{
  basicstyle=\ttfamily\small,
  breaklines=true,
  frame=single,
  xleftmargin=\parindent,
}

\date{}

\begin{document}

% -----------------------------------------------------------------------
% HEADER
% -----------------------------------------------------------------------

\noindent\textbf{Name:}~\rule{0.35\linewidth}{0.4pt}%
\hspace{0.05\linewidth}%
\textbf{Student ID\#:}~\rule{0.20\linewidth}{0.4pt}

\medskip
\noindent\textbf{[COURSE CODE] -- [TERM]: [Exam Name] ([N] pts)}

\ifanswers
  \begin{center}\large\textbf{*** ANSWER KEY --- NOT FOR DISTRIBUTION ***}\end{center}
\fi

% -----------------------------------------------------------------------
% DIRECTIONS
% -----------------------------------------------------------------------

\paragraph*{Directions.}
\textit{Write on the exam paper.
You are allowed one 3''$\times$5'' card of notes (front and back, handwritten only).
\textbf{Absolutely no electronic devices of any kind are allowed on your desk
during the exam.}}

% -----------------------------------------------------------------------
% MULTIPLE CHOICE
% -----------------------------------------------------------------------

\paragraph*{Multiple Choice ([N] pts each).}
\textit{Circle the best answer for each question.
\textbf{Do not} write the answer in the space---it will be marked incorrect.}

\begin{enumerate}

% --- mc: standard 4-option ---
\item Which of the following best describes [concept]?
\ifanswers\quad\textbf{Answer: c}\fi
\begin{enumerate}[label=\alph*.]
  \item [plausible distractor]
  \item [plausible distractor]
  \item [correct answer]
  \item [plausible distractor]
\end{enumerate}

% --- tf: true/false (2 options only) ---
\item [Statement about a concept that is either true or false.]
\ifanswers\quad\textbf{Answer: a (True)}\fi
\begin{enumerate}[label=\alph*.]
  \item True
  \item False
\end{enumerate}

% --- code: code interpretation T/F ---
% Show a complete, self-contained implementation (~15 lines max).
% Incorrect implementations must have a subtle, realistic bug.
% Do not place code questions back-to-back.
\item The following implementation is correct.
\ifanswers\quad\textbf{Answer: b --- [describe the specific bug]}\fi
\begin{lstlisting}
// Show a complete, self-contained implementation.
// Correct or subtly broken -- make the bug realistic.
void example(void)
{
    // ...
}
\end{lstlisting}
\begin{enumerate}[label=\alph*.]
  \item True
  \item False
\end{enumerate}

\newpage

\end{enumerate}

% -----------------------------------------------------------------------
% SHORT ESSAY
% -----------------------------------------------------------------------

\paragraph*{Short essay and code interpretation ([N] pts each).}
\textit{In the space provided below, answer the following questions completely,
but as briefly as possible.
Partial credit will be given for incomplete answers.}

\begin{enumerate}[resume]

\item [Essay Q1: compare/contrast, explain a tradeoff, or analyze a scenario.
       Should require synthesis --- not answerable from a single definition.]
\ifanswers
  \par\medskip\noindent\textbf{Model answer:}
  [Write a complete answer usable as a grading rubric.]
  \par Full credit requires [specific criterion A] AND [specific criterion B].
  \par Partial credit for [partial criterion].
\else
  \bigskip\bigskip\bigskip\bigskip\bigskip
  \bigskip\bigskip\bigskip\bigskip\bigskip
  \bigskip\bigskip\bigskip\bigskip
\fi

\item [Essay Q2: second synthesis or code-interpretation question.]
\ifanswers
  \par\medskip\noindent\textbf{Model answer:}
  [Write a complete answer usable as a grading rubric.]
  \par Full credit requires [specific criterion A] AND [specific criterion B].
  \par Partial credit for [partial criterion].
\else
  \bigskip\bigskip\bigskip\bigskip\bigskip
  \bigskip\bigskip\bigskip\bigskip\bigskip
  \bigskip\bigskip\bigskip\bigskip
\fi

\end{enumerate}

\end{document}
